% Compile with: uplatex file.tex && dvipdfmx file.dvi
%           or: latexmk -uplatex file.tex
\documentclass[dvipdfmx]{article}
\usepackage{amsmath, amssymb}
\usepackage{tikz}

\title{scenario\_spin\_chain\_spectrum}
\date{\today}

\begin{document}
\maketitle

For the Ising model (2 sites), compute the commutator $[H,S^x_1]$ of the spin operator $S^x_1$ — the first step of finding how $S^x$ evolves in time (its equation of motion).

Find the commutator of $S^x_1$ for the Ising model with 2 sites

\begin{equation}
  [H, S^{x}_{1}] = i J S^{y}_{1} S^{z}_{2}
  \label{eq:E1}
\end{equation}

Compute the double commutator $[H,[H,S^x_1]]$ of the same operator (the first-order result feeds directly into the second, corresponding to an "acceleration").

Find the double commutator of $S^x_1$ for the Ising model with 2 sites

\begin{equation}
  [H, [H, S^{x}_{1}]] = J^{2} S^{x}_{1} \left(S^{z}_{2}\right)^{2}
  \label{eq:E2}
\end{equation}

Find the infinite-temperature moments $M_0=\langle(S^x)^2\rangle,\,M_2$, in preparation for expanding the operator's time evolution in a Krylov subspace.

Find the moments up to 2nd order of $S^x_1$ for the Ising model with 2 sites

\begin{equation}
  M_{0} = \frac{1}{4},\ M_{2} = \frac{1}{16} J^{2}
  \label{eq:E3}
\end{equation}

From the moments, find the continued-fraction coefficient $b_1^2$ representing the Green's function of the dynamical correlation function (the recursion method).

Find the continued fraction of $S^x_1$ for the Ising model with 2 sites

\begin{equation}
  b_{1}^2 = \frac{1}{4} J^{2}
  \label{eq:E4}
\end{equation}

From the continued fraction, find the spectral function $S(\omega)$. Since this is a finite system, it should be a sum of discrete delta functions.

Find the spectral function of $S^x_1$ for the Ising model with 2 sites

\begin{equation}
  G(z) = \cfrac{1}{z - \cfrac{\frac{1}{4} J^{2}}{z}},\quad S(\omega) = \frac{1}{2}\,\delta(\omega - \frac{1}{2} J) + \frac{1}{2}\,\delta(\omega + \frac{1}{2} J)
  \label{eq:E5}
\end{equation}

Now find the spectral function at the central site of a 3-site Heisenberg chain. There should be two coefficients this time, giving a central peak plus side peaks.

Find the spectral function of $S^x_2$ for the Heisenberg model on a 3-site chain

\begin{equation}
  G(z) = \cfrac{1}{z - \cfrac{J^{2}}{z - \cfrac{\frac{5}{4} J^{2}}{z}}},\quad S(\omega) = \frac{5}{9}\,\delta(\omega) + \frac{2}{9}\,\delta(\omega - \frac{3}{2} J) + \frac{2}{9}\,\delta(\omega + \frac{3}{2} J)
  \label{eq:E6}
\end{equation}

Without specifying a lattice shape, find the second moment of the Heisenberg model as a general formula in the coordination number $z$.

Find the 2nd moment of $S^x$ for the Heisenberg model on a general lattice

\begin{equation}
  M_0 = \frac{1}{4},\ M_2 = \frac{1}{8} J^{2} z
  \label{eq:E7}
\end{equation}

Find the high-temperature expansion of the Heisenberg model's energy, up to 4th order, as a general formula (in graph invariants such as edge count and triangle count) without specifying a lattice shape.

Find the high-temperature expansion of the Heisenberg model in general form up to 4th order

\begin{equation}
  \langle H^2 \rangle = \frac{3}{16} J^{2} |E|,\ \langle H^3 \rangle = \frac{9}{32} J^{3} N_\triangle - \frac{3}{32} J^{3} |E|,\ \langle H^4 \rangle = \frac{9}{32} J^{4} N_\square - \frac{9}{32} J^{4} N_\triangle - \frac{3}{32} J^{4} N_\wedge - \frac{3}{128} J^{4} |E| + \frac{27}{256} J^{4} |E|^{2}
  \label{eq:E8}
\end{equation}

Apply the general formula to a concrete $3\times3$ square lattice and check consistency with the general formula from the previous step.

Find the high-temperature expansion of the Heisenberg model on a 3x3 square lattice

\begin{equation}
  \langle H^2 \rangle = \frac{9}{4} J^{2},\ \langle H^3 \rangle = -\frac{9}{8} J^{3} \quad (|E| = 12,\ N_\triangle = 0)
  \label{eq:E9}
\end{equation}

Directly count the number of triangles in a 3-vertex graph (a direct check backing up the triangle count $N_\triangle$ appearing in the 3rd-order term of the high-temperature expansion).

Find the number of triangles of the graph (1,2),(2,3),(1,3)

\begin{equation}
  \#\triangle = 1
  \label{eq:E10}
\end{equation}

To contrast with the discrete spectrum (sum of delta functions) of a finite cluster, find the continuous spectral density of an infinite system (the square-root terminator).

Find the continuous spectral density of the square-root terminator

\begin{equation}
  t(z) = \frac{z - \sqrt{z^2 - 4 b_\infty^2}}{2},\quad S(\omega) = \frac{\sqrt{4 b_\infty^2 - \omega^2}}{2\pi b_\infty^2}\quad (|\omega| \le 2 b_\infty)
  \label{eq:E11}
\end{equation}

\end{document}